% =====================================================================
\section{Características del entorno urbano accesible}
\label{sec:caracteristicas}
% =====================================================================

Una vez identificada la norma Orden TMA/851/2021 como la normativa legal de referencia, en esta sección se va a realizar la extracción de los requisitos técnicos específicos para considerar una vía como accesible: itinerarios peatonales accesibles, obras, mobiliario o vados y pasos peatonales.
\\

En cada una de las categorías de características, se describen los parámetros mínimos establecidos, cuáles son los incumplimientos habituales y de qué manera afectan a la diversidad funcional, y finalmente cómo aparecen en los conjuntos de datos de Albacete y Algeciras.
\\

Se identifican también los parámetros que se pueden extraer automáticamente desde imágenes aéreas o transversales, sirviendo de base para las hipótesis de investigación.
\\

% ---------------------------------------------------------------------
\subsection{Itinerario peatonal accesible (IPA)}
\label{subsec:ipa}

El itinerario peatonal accesible (IPA) es la parte de uso peatonal de la vía urbana que garantiza la movilidad segura y circulación, de manera autónoma por parte de todas las personas, independientemente de su diversidad funcional.
\\

\textbf{La anchura libre de paso} es el ancho del área de itinerario peatonal que permite la circulación de peatones totalmente libre de obstáculos, tanto permanentes como temporales.
\\

\textbf{La altura libre de paso} es el alto del trayecto peatonal durante el que las personas pueden circular sin riesgo de impacto contra elementos situados superiormente.
\\

\textbf{La pendiente longitudinal} es la inclinación del itinerario, en porcentaje, en dirección de la marcha.
\\

Mientras que \textbf{la pendiente transversal} es la inclinación del itinerario, en porcentaje, perpendicular a la dirección de la marcha.
\\

% (2) Parámetros que exige la Orden
\begin{tablarequisitos}{Requisitos del itinerario peatonal accesible.}{tab:req-ipa}
  Anchura libre de paso    & Anchura mínima de 1,80 metros. & art. 5.2.a y 5.2.b \\
  \cmidrule(lr){1-3}
  Altura libre de paso     & Altura mínima de 2,20 metros. & art. 5.2.c \\
  \cmidrule(lr){1-3}
  Pendiente longitudinal   & Máxima del 6\%. & art. 5.2.g \\
  \cmidrule(lr){1-3}
  Pendiente transversal    & Máxima del 2\%. & art. 5.2.f \\
  \cmidrule(lr){1-3}
  Resaltes y desniveles    & Sin escalones aislados & art. 5.2.d \\
  \cmidrule(lr){1-3}
  Continuidad              & Sin interrupciones ni obstáculos en todo el recorrido. Uniforme horizontal y vertical. & art. 5.2.a \\
  \cmidrule(lr){1-3}
  Ubicación                & Preferentemente junto a la línea de fachada. En caso de no poder cumplirse, instalar una franja guía de pavimento táctil. & art. 5.2.a, 45 y 46 \\
  \cmidrule(lr){1-3}
  Plataforma única         & Incorporar señalización adecuada para separar la zona peatonal de la zona de circulación. & art. 5.3 \\
\end{tablarequisitos}

\autoref{fig:ipa-plano} ilustra gráficamente cómo se miden y aplican estos parámetros en el itinerario peatonal accesible.
\\

A continuación se indican cuáles son los incumplimientos típicos de estas normas, y su efecto:
% (3) Incumplimientos típicos y su efecto
\begin{inconvenientes}
  \item \textbf{Anchura libre inferior a la exigida:} hace inviable la circulación de personas con movilidad reducida, en silla de ruedas o con carritos de bebé o de la compra.
  \item \textbf{Altura libre insuficiente:} riesgo de impacto en las personas altas, y personas con discapacidad visual.
  \item \textbf{Pendiente longitudinal superior al 6\%:} complicación de la circulación para personas usuarias de silla de ruedas, personas mayores y aquellas con problemas del corazón o esfuerzo físico.
  \item \textbf{Pendiente transversal excesiva:} hace que las personas usuarias de silla de ruedas o diferentes tipos de carrito no sigan una línea recta, desviándose involuntariamente hacia la pendiente, dificultando el tránsito de peatones.
  \item \textbf{Falta de continuidad:} interrumpe la circulación de peatones, dificultando el movimiento de personas ciegas o con diversidad cognitiva.
  \item \textbf{Falta de señalización en plataformas únicas:} incrementa el riesgo de atropello para todo tipo de peatones, con mayor incidencia en las personas ciegas o con problemas cognitivos.
\end{inconvenientes}

% (4) Cómo aparece en los datos
\datosplanes{Dimensiones de pendiente (\texttt{PENDIENTE*}) y anchura (\texttt{ANCHURA*}).}
        {Dimensiones de pendiente (\texttt{PUNT\_PEND\_*}) y anchura (\texttt{PUNT\_ANCH\_*}).}

% ---------------------------------------------------------------------
\subsection{Pavimentos y elementos táctiles indicadores}
\label{subsec:pavimentos}

% (1) Qué es y por qué importa
La elección del pavimento, así como su calidad y estado, son vitales para garantizar la seguridad en la accesibilidad de vías urbanas. Se deben optar por superficies uniformes, que no resbalen y que ofrezcan una estabilidad para todas las personas, independientemente de su diversidad funcional.
\\

Se regulan los tipos de pavimento táctil según su finalidad, distinguiento entre advertencia y encaminamiento. En cuanto a los pavimentos de advertencia, son necesarios para indicar la cercanía a un peligro, desnivel u obstáculo; o fin de una zona de riesgo, alertando a las personas ciegas o con baja visión.
\\

Por otro lado, los pavimentos de encaminamiento se utilizan para orientar y dirigir a las personas, utilizando franjas de inicio y acabado.
\\

% (2) Parámetros que exige la Orden
\begin{tablarequisitos}{Requisitos de pavimentos y elementos táctiles indicadores.}{tab:req-pavimentos}
  Resbaladicidad                     & El material debe cumplir con los requisitos de accesibilidad Clase 3 (resistencia al deslizamiento mayor a 45) del Código Técnico de la Edificación. & art. 11.1, Real Decreto 314/2006, de 17 de marzo \\
  \cmidrule(lr){1-3}
  Continuidad y resaltes             & Superficie dura y estable y sin resaltes superiores a 4 mm. & art. 11.1, 12.1 \\
  \cmidrule(lr){1-3}
  Pavimento táctil de encaminamiento & Piezas de 4 mm de altura y 40 cm de anchura, de relieve lineal. & art. 45.2.a \\
  \cmidrule(lr){1-3}
  Pavimento táctil de advertencia    & Piezas de 4 mm de altura y 40 cm de anchura, de relieve de botones o franjas. & art. 45.2.b \\
\end{tablarequisitos}

% (3) Incumplimientos típicos y su efecto
En cuanto a los incumplimientos:
\begin{inconvenientes}
  \item \textbf{Pavimentos resbaladizos:} aumento de riesgo de caídas, disminución de la movilidad, afectando a personas mayores, gente usuaria de sillas de ruedas o muletas, incluso a toda la población.
  \item \textbf{Resaltes elevados (alcantarillas, baldosines):} riesgo de tropiezos y caídas, afectando principalmente a personas mayores o con problemas de movilidad, también a toda la población - por ejemplo, personas mirando el móvil.
  \item \textbf{Ausencia de pavimento táctil:} en el caso de advertencias, dificultad para detectar peligros, especialmente para personas ciegas; o en el caso de encaminamiento, posibilidad de pérdida de la trayectoria en línea recta, con riesgo de invasión de calzada o cruces con otras personas.
\end{inconvenientes}

% (4) Cómo aparece en los datos
\datosplanes{Dimensión de pavimento (\texttt{PAVIMENTO*}).}
        {Dimensión de pavimento (\texttt{PUNT\_PAVI\_*}).}


% ---------------------------------------------------------------------
\subsection{Vados peatonales, pasos de peatones e isletas de refugio}
\label{subsec:vados}

% (1) Qué es y por qué importa
Los vados, pasos o isletas constituyen los cruces peatonales con el tráfico rodante, y constituyen puntos críticos con alta siniestralidad y mayor riesgo para los peatones si no se diseñan de manera correcta.
\\

En primer lugar, un \textbf{vado peatonal} es el área de la acera destinada al inicio del cruce peatonal, habilitando un acceso seguro a la calzada, generalmente utilizando una inclinación suave para salvar el habitual desnivel entre la acera y la calzada.
\\

El siguiente elemento que se encuentran las personas es el \textbf{paso de peatones}, siendo el área de la calzada destinada al cruce seguro de peatones, suficientemente señalizada y delimitada para garantizar la visibilidad, y en ocasiones ayudada por semaforización.
\\

Por último, las \textbf{isletas de refugio} son áreas intermedias situadas en la calzada habilitadas para permitir a la ciudadanía cruzar la calzada en los tiempos necesarios, sirviendo de lugar seguro de espera, especialmente para calzadas anchas y avenidas.
\\

\begin{figure}[htbp]
\centering
\begin{tikzpicture}[tfm, font=\sffamily\scriptsize,
    every node/.append style={execute at begin node={\begin{hyphenrules}{nohyphenation}},
                              execute at end node={\end{hyphenrules}}},
    rotulo/.style={rotulomapa, fill=papel, draw=tinta-suave, line width=\trazofino, inner sep=1.5pt, align=left},
    guia/.style={draw=tinta, line width=\trazofino},
    vado/.style={fill=seg-acera, draw=tinta-suave, line width=\trazofino,
               postaction={pattern=north east lines, pattern color=tinta-suave}},
    semaforo/.style={circle, fill=seg-obstaculo, draw=papel, line width=0.6pt, inner sep=1.6pt},
  ]
 
% =============================== a) PLANTA ===============================
\node[anchor=south west, font=\sffamily\small\bfseries, text=tinta] at (-0.6,8.45) {a) Planta};
\begin{scope}[x=0.34cm, y=0.34cm]
  % Bandas: edificación, acera, calzada (dos sentidos), isleta, acera, edificación
  \fill[seg-edificio] (0,0)    rectangle (22,0.6);
  \fill[seg-acera]    (0,0.6)  rectangle (22,4.1);
  \fill[seg-calzada]  (0,4.1)  rectangle (22,19.1);
  \fill[seg-acera]    (0,19.1) rectangle (22,22.6);
  \fill[seg-edificio] (0,22.6) rectangle (22,23.2);
  \draw[tinta-suave, line width=\trazofino] (0,0) rectangle (22,23.2) (0,0.6) -- (22,0.6) (0,22.6) -- (22,22.6);
  % Isleta elevada, con el paso rebajado a nivel de calzada
  \fill[seg-acera] (0,10.6) rectangle (9,12.6);  \fill[seg-acera] (13,10.6) rectangle (22,12.6);
  \draw[tinta, line width=\trazo] (0,10.6) -- (9,10.6) -- (9,12.6) -- (0,12.6)
                                   (22,10.6) -- (13,10.6) -- (13,12.6) -- (22,12.6);
  % Marcas viales de carril
  \foreach \y in {7.35,15.85}{
    \draw[papel, line width=0.7pt, dash pattern=on 5pt off 4pt] (0,\y) -- (8.2,\y) (13.8,\y) -- (22,\y);}
  % Paso de peatones (bandas paralelas al bordillo)
  \foreach \y in {4.6,5.6,...,9.6}   {\fill[seg-paso-peatones] (9,\y) rectangle (13,\y+0.5);}
  \foreach \y in {13.1,14.1,...,18.1}{\fill[seg-paso-peatones] (9,\y) rectangle (13,\y+0.5);}
  % Vados: plano inclinado principal + alas laterales
  \filldraw[vado] (8,4.1) -- (14,4.1) -- (13,2.1) -- (9,2.1) -- cycle;
  \filldraw[vado] (8,19.1) -- (14,19.1) -- (13,21.1) -- (9,21.1) -- cycle;
  \draw[tinta-suave, line width=\trazofino, densely dashed] (9,2.1) -- (9,4.1) (13,2.1) -- (13,4.1)
                                                            (9,21.1) -- (9,19.1) (13,21.1) -- (13,19.1);
  % Bordillos
  \draw[tinta, line width=1.2pt] (0,4.1) -- (22,4.1) (0,19.1) -- (22,19.1);
  % Pavimento táctil: advertencia junto al bordillo y franja-guía desde la fachada
  \advertencia{9}{3.1}{13}{3.9}    \direccional{10.5}{0.6}{11.5}{3.1}
  \advertencia{9}{19.3}{13}{20.1}  \direccional{10.5}{20.1}{11.5}{22.6}
  % Sentido de la pendiente del vado (baja hacia la calzada)
  \draw[-{Stealth[length=4pt]}, tinta, line width=0.6pt] (12.3,2.25) -- (12.3,3.0);
  \draw[-{Stealth[length=4pt]}, tinta, line width=0.6pt] (12.3,20.95) -- (12.3,20.2);
  % Semáforos con pulsador
  \node[semaforo] (s1) at (14.2,3.5) {};
  \node[semaforo] (s2) at (7.8,19.7) {};
 
  % ---- Rótulos (sobre las zonas libres de acera, isleta y calzada)
  \node[rotulo, anchor=west] (r1) at (0.4,1.6) {Franja-guía\\direccional\\(acanaladura)};
  \draw[guia] (r1.east) -- (10.6,1.6);
  \node[rotulo, anchor=east] (r2) at (21.7,1.75) {Vado: plano\\inclinado y alas};
  \draw[guia] (r2.west) -- (13.3,2.7);
  \node[rotulo, anchor=west] (r3) at (0.4,20.85) {Semáforo con\\pulsador accesible};
  \draw[guia] (r3.east) -- (s2.west);
  \node[rotulo, anchor=east] (r4) at (21.7,20.85) {Franja de advertencia\\(botones)};
  \draw[guia] (r4.west) -- (12.8,19.7);
  \node[rotulo, anchor=west] at (0.4,11.6) {Isleta (cruce \textgreater\,14\,m)};
  \node[rotulomapa, text=papel, anchor=west] at (0.4,5.6) {Calzada};
  \node[rotulo, anchor=east] (r5) at (21.7,6.0) {Paso de peatones};
  \draw[guia] (r5.west) -- (13,6.85);
  \node[rotulo, anchor=east] (r6) at (21.7,14.6) {Isleta rebajada\\a nivel de calzada};
  \draw[guia] (r6.west) -- (12.5,12.3);
 
  % ---- Cotas
  % Anchura del vado (abajo) y del paso (arriba)
  \draw[cotaref] (9,2.1) -- (9,-1.1)  (13,2.1) -- (13,-1.1);
  \draw[cota] (9,-0.8) -- (13,-0.8) node[midway, cotatexto, below=1pt] {Vado $\geq$\,1,80\,m};
  \draw[cotaref] (9,21.1) -- (9,24.3)  (13,21.1) -- (13,24.3);
  \draw[cota] (9,24.0) -- (13,24.0) node[midway, cotatexto, above=1pt] {Paso $\geq$ ancho de los vados};
  % Anchura del paso en la isleta
  \draw[cota] (9,11.6) -- (13,11.6) node[midway, cotatexto] {$\geq$ paso};
  % Distancia del pulsador al límite exterior del paso
  \draw[cotaref] (13,2.1) -- (13,0.75)  (14.2,3.25) -- (14.2,0.75);
  \draw[cota] (13,0.95) -- (14.2,0.95);
  \node[cotatexto, anchor=north, inner sep=0.8pt] at (13.6,0.8) {$\leq$\,1,50\,m};
  % Longitud de cruce (isleta obligatoria por encima de 14 m)
  \draw[cotaref] (0,4.1) -- (-1.4,4.1)  (0,19.1) -- (-1.4,19.1);
  \draw[cota] (-1.0,4.1) -- (-1.0,19.1) node[midway, cotatexto, rotate=90]
    {Cruce \textgreater\,14\,m $\Rightarrow$ isleta obligatoria};
  % Escala gráfica
  \begin{scope}[shift={(16.5,-1.25)}]
    \escalagrafica{(0,0)}{5}{5\,m}
  \end{scope}
\end{scope}
 
% ======================= b) SECCIÓN POR EL VADO ========================
\begin{scope}[shift={(9.1,5.45)}, x=1.0cm, y=2.6cm]
  \node[anchor=south west, font=\sffamily\small\bfseries, text=tinta] at (-0.3,1.1) {b) Sección por el vado};
  % edificación, acera, plano inclinado y calzada
  \filldraw[seg-edificio, draw=tinta-suave, line width=\trazofino] (-0.3,0.15) rectangle (0,0.95);
  \filldraw[seg-acera, draw=tinta-suave, line width=\trazofino] (0,-0.12) -- (0,0.15) -- (1.5,0.15) -- (3.5,0.0) -- (3.5,-0.12) -- cycle;
  \fill[seg-calzada] (3.5,-0.12) rectangle (4.3,0.0);
  \draw[tinta, line width=\trazo] (0,0.15) -- (1.5,0.15) -- (3.5,0.0) -- (4.3,0.0);
  % pavimento táctil sobre la superficie (trama, como en planta)
  \fill[papel] (0,0.15) rectangle (1.5,0.19);
  \foreach \x in {0.05,0.15,...,1.5}{\draw[tinta, line width=0.4pt] (\x,0.15) -- (\x,0.19);}
  \fill[papel] (1.5,0.15) -- (2.5,0.075) -- (2.5,0.115) -- (1.5,0.19) -- cycle;
  \foreach \t in {0.05,0.15,...,1.0}{\draw[tinta, line width=0.4pt] ({1.5+\t},{0.15-\t*0.075}) -- ++(0,0.04);}
  \fill[papel] (2.5,0.075) -- (3.3,0.015) -- (3.3,0.055) -- (2.5,0.115) -- cycle;
  \foreach \t in {0.06,0.18,...,0.8}{\fill[tinta] ({2.5+\t},{0.075-\t*0.075+0.02}) circle (0.6pt);}
  % cotas
  \draw[cotaref] (2.5,0.12) -- (2.5,0.5)  (3.3,0.06) -- (3.3,0.5);
  \draw[cota] (2.5,0.44) -- (3.3,0.44) node[midway, cotatexto, above=1pt, fill=none] {Advertencia 0,60--1,20\,m};
  \draw[cotaref] (3.5,0.0) -- (3.5,0.3);
  \draw[cota] (3.3,0.26) -- (3.5,0.26);
  \node[cotatexto, anchor=west, fill=none] at (3.55,0.26) {0,10--0,30\,m};
  \node[rotulomapa, align=left, anchor=north west] at (-0.3,-0.17)
    {Pendiente longitudinal $\leq$\,8\,\% (tramo $\leq$\,3\,m)\\
     o $\leq$\,10\,\% (tramo $\leq$\,2\,m); transversal $\leq$\,2\,\%\\
     Resalte con la calzada \textless\,4\,mm};
  \node[rotulomapa, anchor=west] at (0.05,0.3) {franja-guía};
  \node[rotulomapa, text=tinta-suave, anchor=south east, font=\sffamily\tiny] at (4.3,0.75) {escala vertical exagerada};
\end{scope}
 
% ======================= c) ALZADO DEL PULSADOR ========================
\begin{scope}[shift={(10.0,-0.15)}, x=1.1cm, y=1.1cm]
  \node[anchor=south west, font=\sffamily\small\bfseries, text=tinta] at (-1.0,2.95) {c) Pulsador};
  \filldraw[seg-acera, draw=tinta-suave, line width=\trazofino] (-0.9,-0.12) rectangle (2.2,0);
  \draw[tinta, line width=\trazo] (-0.9,0) -- (2.2,0);
  % rango de altura admisible
  \fill[uoc-cian-suave] (-0.3,0.8) rectangle (0.3,1.2);
  \draw[uoc-azul, line width=0.6pt] (-0.3,0.8) -- (0.3,0.8) (-0.3,1.2) -- (0.3,1.2);
  % báculo, semáforo peatonal y pulsador
  \fill[seg-obstaculo] (-0.04,0) rectangle (0.04,2.6);
  \fill[seg-obstaculo, rounded corners=1pt] (-0.12,2.0) rectangle (0.12,2.6);
  \filldraw[papel, draw=tinta, line width=0.5pt, rounded corners=1pt] (0.04,0.92) rectangle (0.2,1.12);
  % cotas de altura
  \draw[cotaref] (-0.3,0.8) -- (-0.75,0.8)  (-0.3,1.2) -- (-0.75,1.2);
  \draw[cota] (-0.65,0) -- (-0.65,0.8);
  \draw[cota] (-0.65,0.8) -- (-0.65,1.2);
  \node[cotatexto, anchor=east, fill=none] at (-0.7,0.8) {0,80\,m};
  \node[cotatexto, anchor=east, fill=none] at (-0.7,1.2) {1,20\,m};
  \node[rotulomapa, align=left, anchor=west] at (0.4,1.0)
    {Pulsador accesible\\a 0,80--1,20\,m};
  \node[rotulomapa, align=left, anchor=west] at (0.4,2.3)
    {Avisador acústico\\homologado\\(a demanda)};
\end{scope}
 
% =============================== Leyenda ===============================
\begin{scope}[shift={(-0.3,-1.15)}]
  \foreach \c/\t/\x in {seg-edificio/Edificación/0, seg-acera/Acera e isleta/2.2,
                        seg-calzada/Calzada/4.7, seg-paso-peatones/Paso de peatones/6.6} {
    \node[muestra, fill=\c] at (\x+0.15,0) {};
    \node[rotulomapa, anchor=west] at (\x+0.32,0) {\t};}
  \node[muestra, fill=seg-acera, postaction={pattern=north east lines, pattern color=tinta-suave}] at (9.65,0) {};
  \node[rotulomapa, anchor=west] at (9.82,0) {Vado};
  \begin{scope}[shift={(0,-0.5)}, x=0.4cm, y=0.4cm]
    \advertencia{0}{-0.35}{0.75}{0.4}   \node[rotulomapa, anchor=west] at (0.95,0) {Advertencia (botones)};
    \direccional{9.5}{-0.35}{10.25}{0.4} \node[rotulomapa, anchor=west] at (10.45,0) {Direccional (acanaladura)};
    \node[semaforo] at (20.6,0) {};     \node[rotulomapa, anchor=west] at (21.05,0) {Semáforo con pulsador};
  \end{scope}
\end{scope}
\end{tikzpicture}

\caption{Vados, pasos de peatones e isletas de refugio: disposición en planta y parámetros
  mínimos (tabla~\ref{tab:req-vados}).}
\label{fig:vados}
\fuente{elaboración propia a partir de la Orden TMA/851/2021, con asistencia de Claude (Anthropic) mediante instrucciones del autor~\cite{claude}.}
\end{figure}

% (2) Parámetros que exige la Orden
\begin{tablarequisitos}{Requisitos de vados, pasos de peatones e isletas.}{tab:req-vados}
  Dimensiones del vado         & Anchura mínima de 1,80 m, con resalte inferior a 4 mm. & art. 20.2 \\
  \cmidrule(lr){1-3}
  Pendiente del vado           & Pendiente longitudinal máxima del 8\% (tramos hasta 3 m), o 10\% (tramos hasta 2 m) en el plano inclinado y transversal máxima del 2\%. & art. 20.6 \\
  \cmidrule(lr){1-3}
  Dimensiones del cruce        & Anchura igual o mayor que los vados que lo inician y finalizan (mínimo 1,80 m). & art. 21.3 \\
  \cmidrule(lr){1-3}
  Pavimento táctil en el cruce & Franja táctil indicador de advertencia (60 a 120 cm) separada de calzada 10 a 30 cm. Franja-guía direccional (80 a 120 cm) desde fachada al centro de advertencia, perpendicular al tráfico y alineada con el lado opuesto. & art. 20 y 21.4 \\
  \cmidrule(lr){1-3}
  Semaforización accesible     & Pulsador accesible (altura entre 0,80 y 1,20 m) a menos de 1,50 m del límite externo y avisador acústico homologado (activado por demanda o mando). & art. 23 \\
  \cmidrule(lr){1-3}
  Isletas de refugio           & Anchura mínima igual al paso de peatones (mínimo 1,80 m) y nivelada con la calzada o con vados. Son obligatorios para distancias de cruce mayores a 14 m. & art. 22 \\
\end{tablarequisitos}

% (3) Incumplimientos típicos y su efecto
Se detallan a continuación los efectos de no cumplir con los requisitos:
\begin{inconvenientes}
  \item \textbf{Desnivel entre calzada y vado peatonal elevado:} riesgo de vuelco para personas usuarias de silla de ruedas, o dificultad (incluso imposibilidad) de acceso a la calzada.
  \item \textbf{Pendiente del vado excesiva:} peligro para las personas usuarias de silla de ruedasm con movilidad reducida o mayores, dificultando o impidiendo el acceso al cruce.
  \item \textbf{Falta o inoperatividad del aviso acústico en semáforos:} hace más difícil el cruce seguro para personas ciegas o con discapacidad visual.
  \item \textbf{Isleta de refugio con ancho insuficiente:} echa atrás a las personas con problemas de movilidad o agilidad, por no poder utilizar las isletas como refugio seguro para cruzar.
\end{inconvenientes}

% (4) Cómo aparece en los datos
\datosplanes{Dimensión de cruces (\texttt{CRUCES*}) Hay además columnas \texttt{FCRSEMF*} y \texttt{FCRVADO*}, pero son enlaces a Google Drive (fotografías).}
        {Dimensión de cruces (\texttt{PUNT\_CRUC\_*}).}

% ---------------------------------------------------------------------
\subsection{Obras e intervenciones}
\label{subsec:obras}

% (1) Qué es y por qué importa
Las obras constituyen, según la \textit{Guía para una política integral de promoción y gestión} del IMSERSO \cite{sala_alonso_2012}, el mayor obstáculo para la accesibilidad peatonal en la vía pública. Además, dada su circustancia temporal, impide a las administraciones a determinar la accesibilidad real de la vía pública en todo momento.
\\

Por este motivo fundamental, se incluye en la normativa un capítulo exclusivo a obras e intervenciones (capítulo X).
\\

\begin{figure}[htbp]
\centering
\begin{tikzpicture}[tfm, font=\sffamily\scriptsize,
    every node/.append style={execute at begin node={\begin{hyphenrules}{nohyphenation}},
                              execute at end node={\end{hyphenrules}}},
    rotulo/.style={rotulomapa, fill=papel, draw=tinta-suave, line width=\trazofino, inner sep=1.5pt, align=left},
    guia/.style={draw=tinta, line width=\trazofino},
    obra/.style={fill=superficie, draw=tinta-suave, line width=\trazofino,
                 postaction={pattern=north west lines, pattern color=tinta-suave}},
    % valla o protección en alto contraste (bandas tinta / papel)
    valla/.style={draw=tinta, line width=2.4pt,
                  postaction={draw=papel, line width=1.4pt, dash pattern=on 3pt off 3pt}},
    luz/.style={star, star points=6, star point ratio=2, fill=papel, draw=tinta,
                line width=0.5pt, inner sep=0.9pt},
    pasamanos/.style={draw=tinta, line width=0.9pt},
  ]

% =============================== a) PLANTA ===============================
\node[anchor=south west, font=\sffamily\small\bfseries, text=tinta] at (-0.1,5.15) {a) Planta};
\begin{scope}[x=0.4cm, y=0.4cm]
  % Superficies
  \fill[seg-calzada]  (0,0)  rectangle (22,6);
  \fill[seg-acera]    (0,6)  rectangle (22,10);
  \fill[seg-edificio] (0,10) rectangle (22,11);
  \draw[tinta-suave, line width=\trazofino] (0,6) rectangle (22,11) (0,10) -- (22,10);
  \draw[papel, line width=0.7pt, dash pattern=on 5pt off 4pt] (0,0.9) -- (22,0.9);
  \draw[tinta, line width=1.2pt] (0,6) -- (22,6);

  % 39.1 · IPA mantenido junto al andamio y tras la obra
  \draw[zonaipa] (0.25,6.45) rectangle (6.75,8.45);
  \draw[zonaipa] (17.25,6.45) rectangle (21.75,8.45);
  % 39.2 · Andamio: proyección y pies protegidos en alto contraste
  \draw[tinta-suave, line width=\trazofino, densely dashed] (0.4,8.65) rectangle (6.6,10);
  \foreach \x in {0.6,2.5,4.4,6.3}{
    \fill[seg-obstaculo] (\x-0.15,8.75) rectangle (\x+0.15,9.05);
    \draw[valla, line width=1.6pt, dash pattern=on 1.5pt off 1.5pt,
          postaction={draw=papel, line width=0.8pt, dash pattern=on 1.5pt off 1.5pt}]
          (\x-0.25,8.65) rectangle (\x+0.25,9.15);}
  % Zona de obras ocupando toda la acera
  \fill[obra] (7,6) rectangle (17,10);
  \node[rotulo] at (12.6,8.6) {Zona de obras};
  % 39.4 · Delimitación de la zona de obras (vallas en alto contraste)
  \draw[valla] (7,9.75) -- (7,10) (7,8.0) -- (7,6) -- (17,6) -- (17,10);
  % 39.5 · Puerta de acceso: abre hacia el interior, no invade el IPA
  \draw[tinta, line width=0.8pt] (7,8.0) -- (8.75,8.0);
  \draw[tinta, line width=\trazofino, densely dashed] (8.75,8.0) arc[start angle=0, end angle=90, radius=1.75];
  % 39.3 · Itinerario alternativo exterior sobre la calzada
  \draw[zonaipa] (4.6,3.75) rectangle (19.4,5.85);
  \node[rotulo] at (12,4.8) {Itinerario alternativo exterior};
  \draw[valla] (4.6,3.5) -- (19.4,3.5);
  \draw[pasamanos] (4.6,3.72) -- (19.4,3.72);
  % Señalización luminosa
  \foreach \x in {4.6,12,19.4} {\node[luz] at (\x,3.15) {};}

  % ---- Cotas
  \draw[cotaref] (4.6,3.5) -- (4.6,-0.6)  (19.4,3.5) -- (19.4,-0.6);
  \draw[cota] (4.6,-0.35) -- (19.4,-0.35) node[midway, cotatexto] {luces cada $\leq$\,50\,m (fuera de escala)};
  \begin{scope}[shift={(16.5,-1.6)}]
    \escalagrafica{(0,0)}{5}{5\,m}
  \end{scope}
  % ---- Rótulos
  \node[rotulo, anchor=west] (r1) at (0.4,7.45) {IPA mantenido};
  \node[rotulo, anchor=west] (r2) at (0.4,1.95) {Vallas estables y detectables\\con pasamanos continuo};
  \draw[guia] (r2.east) -- (9.5,3.45);
  \node[rotulo, anchor=east] (r3) at (21.7,1.95) {Señalización\\luminosa};
  \draw[guia] (r3.north) -- (19.4,2.95);
  \node[rotulo, anchor=south west] (r4) at (5.6,11.4) {Andamio: elementos cubiertos\\y en alto contraste};
  \draw[guia] (2.5,9.15) |- (r4.west);

\end{scope}

% =================== b) SECCIÓN POR EL ITINERARIO ALTERNATIVO ===================
\begin{scope}[shift={(9.5,3.45)}, x=0.5cm, y=0.7cm]
  \node[anchor=south west, font=\sffamily\small\bfseries, text=tinta] at (0.7,2.65) {b) Sección del itinerario alternativo};
  % calzada, acera ocupada por la obra y edificación
  \filldraw[seg-calzada, draw=tinta-suave, line width=\trazofino] (2.4,-0.15) rectangle (6,0);
  \filldraw[seg-acera, draw=tinta-suave, line width=\trazofino] (6,-0.15) rectangle (10,0.15);
  \filldraw[seg-edificio, draw=tinta-suave, line width=\trazofino] (10,-0.15) rectangle (10.5,2.6);
  \fill[obra] (6.15,0.15) rectangle (10,2.0);
  % valla de la obra (borde de acera) y valla del lado del tráfico
  \draw[valla] (6.1,0.15) -- (6.1,2.0);
  \draw[valla] (3.55,0) -- (3.55,1.05);
  \node[luz] at (3.3,1.12) {};
  % pasamanos a 0,90 m y guía inferior
  \fill[tinta] (3.8,0.9) circle (1.4pt);
  \draw[pasamanos] (3.55,0.9) -- (3.8,0.9);
  \fill[tinta] (3.62,0.05) rectangle (3.95,0.2);
  % itinerario y persona
  \draw[zonaipa] (3.75,0) rectangle (5.95,2.2);
  \fill[uoc-azul] (4.95,1.62) circle (0.11);
  \fill[uoc-azul, rounded corners=2pt] (4.8,0.02) rectangle (5.1,1.45);
  % cotas
  \draw[cotaref] (3.8,0.9) -- (3.15,0.9);
  \draw[cota] (3.3,0) -- (3.3,0.9);
  \node[cotatexto, anchor=east, fill=none, align=right] at (3.2,0.45) {pasamanos\\0,90\,m};
  \node[rotulomapa, anchor=south east, align=right, inner sep=1pt] at (3.3,1.3) {valla $\geq$\,0,90\,m\\con luz exterior};
  \draw[cotaref] (3.75,0) -- (3.75,-0.55)  (5.95,0) -- (5.95,-0.55);
  \draw[cota] (3.75,-0.45) -- (5.95,-0.45) node[cotatexto, anchor=west, at end, xshift=2pt, fill=none] {$\geq$\,1,80\,m};
  % rótulos
  \node[rotulo] at (8.05,1.1) {Zona de obras};
  \draw[guia] (3.95,0.12) -- (3.2,-0.45) node[rotulomapa, anchor=east, inner sep=1pt] {guía inferior};
\end{scope}

% ======================== c) SECCIÓN POR EL ANDAMIO ========================
\begin{scope}[shift={(9.5,0.35)}, x=0.5cm, y=0.7cm]
  \node[anchor=south west, font=\sffamily\small\bfseries, text=tinta] at (0.7,2.65) {c) Sección por el andamio};
  \filldraw[seg-calzada, draw=tinta-suave, line width=\trazofino] (2.4,-0.15) rectangle (6,0);
  \filldraw[seg-acera, draw=tinta-suave, line width=\trazofino] (6,-0.15) rectangle (10,0.15);
  \filldraw[seg-edificio, draw=tinta-suave, line width=\trazofino] (10,-0.15) rectangle (10.5,2.75);
  % andamio: pies, plataforma y marquesina sobre el itinerario
  \fill[seg-obstaculo] (8.65,0.15) rectangle (8.8,2.6);
  \fill[seg-obstaculo] (9.85,0.15) rectangle (10,2.6);
  \fill[seg-obstaculo] (6.6,2.45) rectangle (10,2.55);
  \draw[seg-obstaculo, line width=0.8pt] (8.72,2.5) -- (9.92,1.4);
  % protección en alto contraste hasta 2,20 m
  \draw[valla, line width=2.8pt, postaction={draw=papel, line width=1.6pt, dash pattern=on 3pt off 3pt}]
       (8.72,0.15) -- (8.72,2.2);
  % IPA bajo la marquesina
  \draw[zonaipa] (6.45,0.15) rectangle (8.5,2.2);
  \fill[uoc-azul] (7.5,1.77) circle (0.11);
  \fill[uoc-azul, rounded corners=2pt] (7.35,0.17) rectangle (7.65,1.6);
  % cotas
  \draw[cotaref] (6.45,2.2) -- (5.7,2.2);
  \draw[cota] (5.85,0.15) -- (5.85,2.2) node[midway, cotatexto, rotate=90] {$\geq$\,2,20\,m};
  \draw[guia] (8.72,1.3) -- (7.1,-0.55) node[rotulomapa, anchor=north, align=center, inner sep=1pt]
    {pie protegido en alto contraste\\hasta 2,20\,m};
\end{scope}

% =============================== Leyenda ===============================
\begin{scope}[shift={(0.1,-1.45)}]
  \foreach \c/\t/\x in {seg-edificio/Edificación/0, seg-acera/Acera/2.2, seg-calzada/Calzada/3.8,
                        seg-obstaculo/Andamio/5.7} {
    \node[muestra, fill=\c] at (\x+0.15,0) {};
    \node[rotulomapa, anchor=west] at (\x+0.32,0) {\t};}
  \node[muestra, obra] at (7.65,0) {};  \node[rotulomapa, anchor=west] at (7.82,0) {Zona de obras};
  \draw[zonaipa] (10.0,-0.12) rectangle (10.3,0.12); \node[rotulomapa, anchor=west] at (10.4,0) {Itinerario accesible};
  \begin{scope}[shift={(0,-0.5)}]
    \draw[valla] (0,0) -- (0.5,0);           \node[rotulomapa, anchor=west] at (0.6,0) {Valla o protección en alto contraste};
    \draw[pasamanos] (5.6,0) -- (6.1,0);     \node[rotulomapa, anchor=west] at (6.2,0) {Pasamanos};
    \node[luz] at (8.25,0) {};               \node[rotulomapa, anchor=west] at (8.45,0) {Señalización luminosa};
    \draw[tinta, line width=0.8pt] (11.6,-0.12) -- (11.9,-0.12);
    \draw[tinta, line width=\trazofino, densely dashed] (11.9,-0.12) arc[start angle=0, end angle=90, radius=0.3];
    \node[rotulomapa, anchor=west] at (12.0,0) {Puerta};
  \end{scope}
\end{scope}
\end{tikzpicture}

\caption{Obras e intervenciones en la vía pública: continuidad del itinerario peatonal
  accesible, itinerario alternativo, delimitación y protección (tabla~\ref{tab:req-obras}).}
\label{fig:obras}
\fuente{elaboración propia a partir de la Orden TMA/851/2021, con asistencia de Claude (Anthropic) mediante instrucciones del autor~\cite{claude}.}
\end{figure}

% (2) Parámetros que exige la Orden
\begin{tablarequisitos}{Requisitos de las obras e intervenciones en la vía pública.}{tab:req-obras}
  Itinerario peatonal accesible   & O mantenerlo accesible o demarcar uno alternativo señalizado según lo indicado en art. 5. & art. 39.1 \\
  \cmidrule(lr){1-3}
  Andamios y estructuras provisionales & Elementos externos cubiertos frente a golpes, garantizando accesibildad con colores alto contraste. & art. 39.2 \\
  \cmidrule(lr){1-3}
  Itinerario alternativo exterior  & Pasamanos continuo a 90 cm altura y guía inferior; o franja-guía de pavimento táctil indicador. & art. 39.3 \\
  \cmidrule(lr){1-3}
  Delimitación de zonas de obras   & Vallas estables y detectables, altura mín. 90 cm. Señalización luminosa cada 50 m, con color contrastado. & art. 39.4 \\
  \cmidrule(lr){1-3}
  Puertas, portones y accesos      & No invadirán el itinerario peatonal accesible. Elementos sobresalientes protegidos con color contrastado hasta 2,20 m de altura. & art. 39.5 \\
\end{tablarequisitos}

% (3) Incumplimientos típicos y su efecto
Los problemas que pueden ocasionar los incumplimientos son:
\begin{inconvenientes}
  \item \textbf{Ancho del itinerario provisional (provisional) insuficiente:} mismos problemas que en un IPA de menos de 1,80 m.
  \item \textbf{Falta de elementos táctiles indicadores:} pueden provocar caídas o desorientación en personas con discapacidad visual o con problemas cognitivos.
  \item \textbf{Elementos de protección mal cubiertos:} pueden provocar accidentes o caídas de elementos que impacten en peatones, de todo tipo de diversidad.
\end{inconvenientes}

% (4) Cómo aparece en los datos (aquí: tratamiento en el análisis)
\datosplanes{No existe una dimensión de obras. Solo aparece en algunas calles que están de obras en los campos \texttt{TIPOLACT} o \texttt{PROPACT}.}
        {No existe una dimensión de obras. Solo hay una calle que se indica que está en obras (en el campo\texttt{PROPUESTA}).}

% ---------------------------------------------------------------------
\subsection{Mobiliario urbano y terrazas}
\label{subsec:mobiliario}

% (1) Qué es y por qué importa
El mobiliario urbano y las terrazas de hostelería consituyen un punto crítico en la accesibilidad de las calles, al ubicarse en el ancho de las aceras.
\\

Por lo que su disposición y delimitaciones se deben regular cuidadosamente para no obstaculizar las trayectorias peatonales accesibles.
\\

% (2) Parámetros que exige la Orden
\begin{tablarequisitos}{Requisitos del mobiliario urbano y las terrazas.}{tab:req-mobiliario}
  Ubicación del mobiliario & Alineado en una banda específica fuera del itinerario peatonal accesible, preferentemente colindante al bordillo. & art. 25.1 \\
  \cmidrule(lr){1-3}
  Dimensiones y alturas    & Elementos elevados deben tener una altura libre $\ge$ 2,20 m o estar protegidos hasta el suelo detectables con bastón. & art. 25.2 \\
  \cmidrule(lr){1-3}
  Franja libre de paso     & Se debe mantener la anchura mínima del IPA (1,80 m), respetando la línea de fachada. & art. 25.1 y art. 5 \\
  \cmidrule(lr){1-3}
  Ocupación por terrazas   & Terrazas, mamparas o puestos temporales no podrán invadir el IPA, ni ocultar la señalización. & art. 25.3 \\
  \cmidrule(lr){1-3}
  Bolardos                 & Alineados, con altura entre 0,75 y 1,00 m y color contrastado diferenciador. & art. 29 \\
\end{tablarequisitos}

% (3) Incumplimientos típicos y su efecto
Se listan a continuación los incumplimientos y su efecto:
\begin{inconvenientes}
  \item \textbf{Mobiliario no alineado:} elimina previsibilidad en el trayecto, y dificulta la movilidad a personas con baja visión o con diversidad cognitiva.
  \item \textbf{Invasión de la franja peatonal por terrazas:} reduce la anchura libre establecida en el artículo 5, con los mismos problemas indicados anteriormente.
  \item \textbf{Elementos elevados invadiendo la franja peatonal:} peligro de impacto con los peatones, especialmente en las personas altas y con discapacidad visual.
\end{inconvenientes}

% (4) Cómo aparece en los datos
\datosplanes{Dimensión de mobiliario (\texttt{MOBILIAR*}). Las columnas \texttt{FMOB*} son enlaces a Google Drive de fotos o informes.}
        {Dimensión de mobiliario (\texttt{PUNT\_MOBI\_*}).}

% ---------------------------------------------------------------------
\subsection{Calles de plataforma única}
\label{subsec:plataforma}

% (1) Qué es y por qué importa
Las plataformas únicas son una solución para las calles estrechas cuyo ancho impida la construcción de un itinerario peatonal accesible con bordillos, en base a los mínimos requisitos de la normativa.
\\

Necesitan combinar la integración de la accesibilidad de las calles con la convivencia entre peatones y vehículos. Tienen prioridad peatonal, y la acera y la calzada se encuentran al mismo nivel, para evitar barreras físicas que penalicen al tránsito peatonal.
\\

La guía de la ONCE distingue dos tipos \cite{once_2024}: \textbf{tránsito compartido}, donde toda la superficie es de prioridad peatonal con una zona delimitada para los vehículos; y \textbf{tránsito diferenciado}, donde acera y calzada se dividen utilizando pavimentos diferentes, con franjas de aviso o mobiliario separador.
\\

% (2) Parámetros que exige la Orden

\textbf{Plataforma única de tránsito compartido}

\begin{tablarequisitos}{Requisitos de plataforma única de tránsito compartido.}{tab:req-plataforma-compartido}
  Velocidad máxima                & Como mucho 20 km/h. & art. 5.3, ONCE 2024 \\
  \cmidrule(lr){1-3}
  Prioridad peatonal y organización & Las personas pueden caminar y trasladarse por toda la zona. Solo se autoriza el paso de un grupo restringido de vehículos (residentes, carga y descarga o emergencias; pero nunca transporte urbano). & ONCE 2024, 5.1--5.4 \\
  \cmidrule(lr){1-3}
  Pavimento homogéneo             & Se usa el mismo pavimento en toda la plataforma; el mobiliario debe estar alineado uniformemente. & ONCE 2024, 5.5 \\
  \cmidrule(lr){1-3}
  Zona de entrada y salida        & Es la zona que necesita más seguridad. Anchura mínima de 1,80 m desde fachada a ambos lados. Banda tacto-visual de advertencia de 40 a 60 cm o banda de encaminamiento. & ONCE 2024, 5.6 \\
\end{tablarequisitos}

\medskip
\textbf{Plataforma única de tránsito diferenciado}

\begin{tablarequisitos}{Requisitos de plataforma única de tránsito diferenciado.}{tab:req-plataforma-diferenciado}
  Pavimento                 & Acera y calzada diferenciadas por pavimento distinto (color y textura). & ONCE 2024, 6.1 \\
  \cmidrule(lr){1-3}
  Bandas delimitación       & Banda tacto-visual de advertencia de 40 a 60 cm. Mobiliario alineado con separación máxima 120 cm. & art. 46, ONCE 2024, 6.2 \\
  \cmidrule(lr){1-3}
  Señalización normalizada  & Señalización vertical y horizontal indicando velocidad máxima, estacionamiento autorizado y delimitación peatón-vehículo. & art. 5.3, ONCE 2024, 6.4 \\
\end{tablarequisitos}

% (3) Incumplimientos típicos y su efecto
En cuanto a los incumplimientos típicos, se encuentran:
\begin{inconvenientes}
  \item \textbf{Vehículos no respetan la prioridad:} hace añicos la plataforma única, y transmite inseguridad y un elevado riesgo de accidentes a todas las personas, afectando especialmente a las personas ciegas.
  \item \textbf{Ausencia de bandas de delimitación en tránsito diferenciado:} potencial invasión de la zona delimitada a los vehículos por parte de peatones, con riesgo de atropello, afectando principalmente a las personas con baja visión.
\end{inconvenientes}

% (4) Cómo aparece en los datos
\datosplanes{No es una dimensión de puntuación, sino una categoría de \texttt{TIPOLACT} («Plataforma»).}
        {Solo existe con el valor «Plataforma única» de \texttt{TIPOLOGÍA}.}

% ---------------------------------------------------------------------
\subsection{Señalización y elementos de información}
\label{subsec:senalizacion}

% (1) Qué es y por qué importa
La normativa estatal establece requisitos específicos para la señalización y los carteles informativos, para que puedan ser detectados y leídos por todas las personas, independientemente de su diversidad. Cabe recordar que la accesibilidad universal necesita tener en cuenta no solo los problemas de movilidad, sino aspectos auditivos, para las personas ciegas y con baja audición o aspectos de adaptación de la lectura, para personas ciegas.
\\
% (2) Parámetros que exige la Orden
\begin{tablarequisitos}{Requisitos de señalización y elementos de información.}{tab:req-senalizacion}
  Señalización visual & Fuentes legibles, tamaño según distancia (0,7 a 7,0 cm), altura entre 0,90 y 1,20 m, con mínimo contraste cromático, evitando reflejos o sombras. & art. 41 \\
  \cmidrule(lr){1-3}
  Información escrita & Priorizar lectura fácil, con textos concisos con pictogramas. Uniformidad de los paneles en el municipio. Nombres de las vías obligatorios en todos los cruces. & art. 40, 41, 42 \\
  \cmidrule(lr){1-3}
  Señalización táctil & Información en Braille y alto relieve en todos los carteles. & art. 44 \\
  \cmidrule(lr){1-3}
  Información sonora & Señalización acústica audible en cruces, acompañada de señalización visual equivalente. & art. 41 \\
  \cmidrule(lr){1-3}
  Elementos transparentes & Obligación de señalizar con bandas opacas de 5 a 10 cm a alturas: entre 0,85 y 1,10 m la primera y de 1,50a 1,70 m la segunda. & art. 42 \\
\end{tablarequisitos}

% (3) Incumplimientos típicos y su efecto
Se detallan los incumplimientos y su efecto en las personas:
\begin{inconvenientes}
  \item \textbf{Falta de paneles informativos o a mala altura:} genera problemas de orientación a todas las personas. Una mala colocación impide que personas en silla de ruedas o con problemas visuales alcancen la información.
  \item \textbf{Falta de contraste cromático y visual:} hace difícil o imposible la lectura a personas con baja visión o daltónicas, por no poder distinguir los textos o pictogramas.
  \item \textbf{Ausencia de información en Braille o altorrelieve:} provoca problemas de orientación a personas ciegas o de baja visión, al no poder leer la información de los paneles.
  \item \textbf{Falta de señalización acústica en cruces:} dificulta el cruce seguro a personas con discapacidad visual, al no saber directamente cuándo es posible iniciar el cruce.
  \item \textbf{Cristales sin señalización táctil:} aumenta el riesgo de colisión contra la superficie transparente para todas las personas, pero especialmente para personas ciegas o con problemas de visión.
\end{inconvenientes}

% (4) Cómo aparece en los datos
\datosplanes{Dimensión de señalización (\texttt{SEÑALIZ*}).}
        {Dimensión de señalización (\texttt{PUNT\_SEÑA\_*}).}

% ---------------------------------------------------------------------
\subsection{Iluminación y visibilidad}
\label{subsec:iluminacion}

% (1) Qué es y por qué importa
La iluminación influye de manera notable en la seguridad de las personas cuando se desplazan por la vía pública en horas de baja visibilidad natural, como las horas nocturnas o los días con niebla, lluvia o cielos cubiertos.
\\

Aunque la orden de referencia TMA/851/2021 establece ciertos requisitos de iluminación para las vías accesibles, las características específicas de los niveles de iluminación están regulados por el Real Decreto 1890/2008 sobre eficiencia energética.
\\

La mayoría de estos parámetros necesitan que los equipos técnicos necesiten desplazarse in situ, para el uso de luxómetros; por lo que son complicados de obtener mediante el uso de imágenes aéreas o a pie de calle.
\\

% (2) Parámetros que exige la Orden
\begin{tablarequisitos}{Requisitos de iluminación en itinerarios peatonales accesibles.}{tab:req-iluminacion}
  Niveles de iluminación & Ajustados a los especificados en \cite{rd_1890_2008} sobre eficiencia energética en alumbrado exterior. & art. 5.2.h \\
  \cmidrule(lr){1-3}
  Uniformidad luminosa & Distribución homogénea de luz sin zonas oscuras, sombras o deslumbramientos que dificulten la visibilidad. & \cite{rd_1890_2008} \\
  \cmidrule(lr){1-3}
  Detección de obstáculos & Iluminación suficiente que permita detectar posibles obstáculos como mobiliario, rejillas, marquesinas o cambios de nivel. & art. 5.2.h \\
\end{tablarequisitos}

% (3) Incumplimientos típicos y su efecto
Se listan los incumplimientos y sus efectos:
\begin{inconvenientes}
  \item \textbf{Niveles de iluminación insuficientes o irregulares:} afecta a todas las personas, cuando los niveles de luz natural no permiten la visión, pudiendo provocar accidentes o caídas.
  \item \textbf{Deslumbramientos o reflejos:} pueden ocasionar molestias así como posibles caídas, especialmente en aquellas personas con baja visión.
  \item \textbf{Falta de iluminación en cruces u obstáculos:} ocasiona dificultad, miedo o impedimento de cruzar las calzadas, así como posible riesgo de accidente o impacto, en todas las personas, pero especialmente en las que tienen problemas de visión.
\end{inconvenientes}

% (4) Cómo aparece en los datos
\datosplanes{Dimensión de iluminación (\texttt{ILUMIN*}).}
        {Dimensión de iluminación (\texttt{PUNT\_ILUM\_*}).}

% =====================================================================
\section{Análisis exploratorio de los datos}
\label{sec:analisis-exploratorio}

Para poder evaluar con qué datos se cuenta para la realización del trabajo, se realiza un análisis exploratorio de los conjuntos de datos aportados por el tutor, de las ciudades de Albacete y Algeciras.

\medskip
\noindent\textit{Origen de los datos.} Los conjuntos de datos han sido facilitados por el tutor, y son los resultados de los análisis de accesibilidad del viario urbano (y de los edificios públicos) realizados por los equipos técnicos contratados por los ayuntamientos.

\medskip
\noindent\textit{Observaciones del análisis.} Al realizar el análisis exploratorio, en ambas ciudades se observaron patrones similares en cuanto a las variables evaluadas, y las variables de resultados o clase. Por esta similitud, se ha inferido que:
\begin{itemize}
  \item o bien existe una metodología común para realizar estos análisis de accesibilidad,
  \item o bien se han realizado por la misma empresa.
\end{itemize}

\medskip
\noindent\textit{Verificación de la empresa adjudicataria.} Tras consultar la Plataforma de Contratación del Estado, se ha podido verificar que en ambas ciudades la empresa adjudicataria de los contratos fue la misma, ACCESIA~\cite{contrato_albacete_2021,contrato_algeciras_2025}. Los expedientes se recogen en la Tabla~\ref{tab:contratos}:

\begin{table}[H]
  \centering
  \caption{Contratos de los análisis de accesibilidad de Albacete y Algeciras.}
  \label{tab:contratos}
  \small
  \renewcommand{\arraystretch}{1.15}
  \begin{tabularx}{\textwidth}{l c >{\raggedright\arraybackslash}X}
    \toprule
    \cabeceratabla
    \textbf{Ciudad} & \textbf{Año del contrato} & \textbf{Expediente} \\
    \midrule
    Albacete  & 2021 & \href{https://contrataciondelestado.es/wps/poc?uri=deeplink:detalle_licitacion\&idEvl=B\%2BdESdnRNPBvYnTkQN0\%2FZA\%3D\%3D}{Expediente de Albacete} \\
    Algeciras & 2025 & \href{https://contrataciondelestado.es/wps/poc?uri=deeplink:detalle_licitacion\&idEvl=s4jpjymDFxh\%2FR5QFTlaM4A\%3D\%3D}{Expediente de Algeciras} \\
    \bottomrule
  \end{tabularx}
  \fuente{Plataforma de Contratación del Estado~\cite{contrato_albacete_2021,contrato_algeciras_2025}.}
\end{table}

\subsection{Resultados del análisis exploratorio}

Los conjuntos de datos contienen una fila por cada calle o tramo de calle, con una puntuación de accesibilidad en cada \gls{dimension-accesibilidad} de la \gls{orden-tma-851-2021}~\cite{mitma_orden_tma851_2021}, incluyendo anchura, pendiente, iluminación, pavimento, cruces, señalización, mobiliario y transporte (únicamente Albacete); y en cada dimensión, evaluando diferentes tipos de \gls{diversidad-funcional}.

\medskip
\noindent\textit{Tipos de diversidad funcional.}
\begin{itemize}
  \item \textbf{Algeciras:} las tipologías de diversidad evaluadas incluyen 3 tipos diferentes: funcional, sensorial y cognitiva.
  \item \textbf{Albacete:} hay 5 tipos diferentes: silla de ruedas, movilidad reducida, visión, hipoacusia y cognitiva.
\end{itemize}

\noindent\textit{Escala de puntuación.}
\begin{itemize}
  \item \textbf{Algeciras:} escala de 0 a 5, en orden ascendente, siendo 0 la peor accesibilidad y 5 la mejor.
  \item \textbf{Albacete:} escala de 0 a 100, siguiendo un orden también ascendente.
\end{itemize}
En ambos casos, las puntuaciones por cada dimensión y tipo de diversidad funcional son números enteros, por lo que se ha inferido que existe un algoritmo que ha utilizado la empresa para ir dando o restando puntuación según se cumple o no con los requisitos de la normativa legal.

\medskip
\noindent\textit{Unidad de cada fila.} En ambos conjuntos se incluye una variable con el nombre de la calle, aportando la puntuación por calle, con algunas calles separadas por tramos. Algeciras también incluye el nombre de la barriada, y algunas calles están separadas por barriadas.

\medskip
\noindent\textit{Tipología de la vía.} Cada calle cuenta con una columna que indica la tipología de la vía, es decir, si se trata de:
\begin{itemize}
  \item una calle peatonal,
  \item de \gls{plataforma-unica},
  \item o de tráfico diferenciado (Algeciras distingue entre calles con un carril, o con más, pero no Albacete),
  \item así como otros tipos como plaza, bulevar, o en obras.
\end{itemize}

\medskip
\noindent\textit{Presupuesto y propuesta de actuación.} En ambos casos se indica en una variable el presupuesto en euros asignado a cada calle tras la evaluación de accesibilidad; y una propuesta de actuación, en la que se distingue entre:
\begin{itemize}
  \item mejoras en la vía existente,
  \item cambio de tipología,
  \item o no realizar ninguna actuación.
\end{itemize}



Para trabajar con ambas ciudades a la vez, se unifican sus columnas en un único conjunto de datos, cuya correspondencia se muestra en la Tabla~\ref{tab:columnas-unificadas}.

\begin{table}[H]
  \centering
  \caption{Columnas del conjunto de datos unificado de Albacete y Algeciras.}
  \label{tab:columnas-unificadas}
  \small
  \renewcommand{\arraystretch}{1.1}
  \setlength{\extrarowheight}{0pt}
  \begin{tabularx}{\textwidth}{>{\raggedright\arraybackslash}p{3.9cm} >{\raggedright\arraybackslash}p{3.6cm} >{\raggedright\arraybackslash}X >{\raggedright\arraybackslash}X}
    \toprule
    \cabeceratabla
    \textbf{Concepto} & \textbf{Columna unificada} & \textbf{Algeciras} & \textbf{Albacete} \\
    \midrule
    \multicolumn{4}{l}{\sffamily\footnotesize\bfseries\color{uoc-azul}Identificación y localización} \\
    Ciudad (nueva, para unir los datos) & \texttt{ciudad} & -- & -- \\
    \cmidrule(lr){1-4}
    Nombre de la calle & \texttt{nombre\_calle} & \texttt{NOMBRE} & \texttt{NOMBRE} \\
    \cmidrule(lr){1-4}
    Nombre de la calle sin la marca de tramo & \texttt{nombre\_calle\_base} & \texttt{NOMBRE} sin la marca de tramo & \texttt{NOMBRE} sin la marca de tramo \\
    \cmidrule(lr){1-4}
    Tramo de la calle (número) & \texttt{tramo} & Letra final de \texttt{REFERENCIA} o marca del nombre & Número del final de \texttt{NOMBRE} \\
    \cmidrule(lr){1-4}
    Barrio & \texttt{barrio} & \texttt{Barriada} & -- (pendiente de cruzar con el callejero) \\
    \cmidrule(lr){1-4}
    Zona & \texttt{zona} & -- & Letras de \texttt{REFER} \\
    \midrule
    \multicolumn{4}{l}{\sffamily\footnotesize\bfseries\color{uoc-azul}Tipología y actuación} \\
    Tipología de la vía & \texttt{tipologia\_via} & \texttt{TIPOLOGÍA} & \texttt{TIPOLACT} \\
    \cmidrule(lr){1-4}
    Propuesta de actuación & \texttt{propuesta\_actuacion} & \texttt{PROPUESTA} & \texttt{PROPACT} \\
    \cmidrule(lr){1-4}
    Presupuesto & \texttt{presupuesto\_euros} & \texttt{PRESUPUESTO} & \texttt{PRESUP} \\
    \midrule
    \multicolumn{4}{l}{\sffamily\footnotesize\bfseries\color{uoc-azul}Puntuaciones y clase} \\
    Puntuación total & \texttt{puntuacion\_total} & \texttt{PUNT\_TOTAL} & \texttt{ACCGLOBAL} \\
    \cmidrule(lr){1-4}
    Puntuación total de un tipo de diversidad & \texttt{puntuacion\_total\_*} & \texttt{PUNT\_TOTAL\_FU} & Media de \texttt{GLOBALSR} y \texttt{GLOBALMR} \\
    \cmidrule(lr){1-4}
    Puntuación global de una dimensión & \texttt{puntuacion\_*\_global} & \texttt{PUNT\_PEND\_GB} & \texttt{GLOBALPEND} \\
    \cmidrule(lr){1-4}
    Puntuación de una dimensión y un tipo de diversidad & \texttt{puntuacion\_*\_*} & \texttt{PUNT\_PEND\_FU} & Media de \texttt{PENDIENTESR} y \texttt{PENDIENTEMR} \\
    \cmidrule(lr){1-4}
    \gls{clase-accesibilidad} total & \texttt{clase\_total} & \texttt{PIC\_TOTAL} & Calculada a partir de \texttt{ACCGLOBAL} \\
    \bottomrule
  \end{tabularx}
  \fuente{elaboración propia a partir de los conjuntos de datos de Albacete y Algeciras.}
\end{table}

\medskip
\noindent\textit{Cómo leer la tabla.}
\begin{itemize}
  \item \textbf{Columnas \texttt{puntuacion\_*\_*}:} el primer comodín es una de las dimensiones (pendiente, anchura, pavimento, cruces, señalización, iluminación, mobiliario y transporte, esta última solo en Albacete) y el segundo uno de los tipos de diversidad (funcional, sensorial y cognitiva).
  \item \textbf{Escala:} las puntuaciones están normalizadas a la escala de 0 a 1.
  \item \textbf{Tipos de diversidad en Albacete:} para unificarlos, funcional es la media de silla de ruedas y movilidad reducida, sensorial la media de visión e hipoacusia, y cognitiva es la de diversidad cognitiva.
  \item \textbf{Transporte:} los ceros de transporte, que indican que no hay datos, se sustituyen por valores nulos.
  \item \textbf{Tramo:} identifica la parte de una calle que está dividida (por ejemplo, la letra B de Algeciras es el tramo 2), y es nulo si la calle no tiene marca de tramo.
\end{itemize}

% =====================================================================
\section{Viabilidad de calcular las dimensiones a partir de imágenes}
\label{sec:viabilidad}
% =====================================================================
\estadobloque{pendiente}{Borrador para revisión. Conecta con OP3 y fundamenta H1, H2 y H5.}

Para poder realizar la comprobación científica de las hipótesis de investigación, se necesita discernir cuáles de las dimensiones que tienen puntuación en el informe (pendiente, anchura, pavimento, cruces, mobiliario, señalización, iluminación; y transporte en Albacete) pueden obtenerse a partir de imágenes aéreas o a pie de calle, y cuáles necesitan un análisis \textit{in situ}.
\\

La tabla~\ref{tab:viabilidad} indica, para cada una de las dimensiones con puntuación en los datos de Albacete y Algeciras cuáles son los artículos de la Orden TMA/851/2021 \cite{mitma_orden_tma851_2021} que la regulan, y si pueden extraerse potencialmente a partir de las imágenes.
\\

El grado de potencialidad de extracción de imágenes se expresa como un porcentaje aproximado de lo que se podría obtener de cada dimensión, y es un valor subjetivo estimado por el autor de este trabajo. Este resultado permitirá el análisis de las hipótesis, obteniendo el conjunto de dimensiones que se utilizarán en caso de imágenes aéreas, a pie de calle o in situ.
\\

Idealmente, el grado se podrá calcular de manera empírica cuando se realicen los modelos de visión computacional de extracción de datos geométricos de las imágenes, que se deja como futura línea de investigación.
\\

\begingroup
\footnotesize
\setlength{\parindent}{0pt}
\renewcommand{\arraystretch}{1.05}
\setlength{\tabcolsep}{3pt}
\begin{longtable}{>{\raggedright\arraybackslash}m{6.2cm} >{\raggedright\arraybackslash}m{1.5cm} >{\raggedright\arraybackslash}m{1.5cm} >{\raggedright\arraybackslash}m{\dimexpr\textwidth-9.2cm-8\tabcolsep\relax}}
  \caption{Viabilidad de obtención de puntuación a partir de las imágenes de cada dimensión de accesibilidad.}\label{tab:viabilidad}\\
  \toprule
  \cabeceratabla
  \textbf{Qué hay que medir} & \textbf{Aérea} & \textbf{Calle} & \textbf{Razón} \\
  \midrule
  \endfirsthead
  \toprule
  \cabeceratabla
  \textbf{Qué hay que medir} & \textbf{Aérea} & \textbf{Calle} & \textbf{Razón} \\
  \midrule
  \endhead
  \bottomrule
  \endfoot
  \bottomrule
  \multicolumn{4}{p{\dimexpr\textwidth-2\tabcolsep\relax}}{\notatabla{Fuente: elaboración propia.}}
  \endlastfoot
    \multicolumn{1}{l}{\sffamily\footnotesize\bfseries\color{uoc-azul}Pendiente} & \textbf{0\,\%} & \textbf{40\,\%} & \\
    \Gls{pendiente-longitudinal} ($\leq$~6\,\%) & \cumple{no}\,0\,\% & \cumple{parcial}\,50\,\% & Desde el aire no se estima; a pie de calle, con varias fotografías podría hacerse una medición geométrica. \\
    \arrayrulecolor{black!25}\cmidrule(lr){1-4}\arrayrulecolor{black}
    \Gls{pendiente-transversal} ($\leq$~2\,\%) & \cumple{no}\,0\,\% & \cumple{parcial}\,50\,\% & Igual que la longitudinal: solo a pie de calle y con varias fotografías. \\
    \arrayrulecolor{black!25}\cmidrule(lr){1-4}\arrayrulecolor{black}
    Pendiente del \glslink{vado-peatonal}{vado} (10\,\% hasta 2\,m, 8\,\% hasta 3\,m) & \cumple{no}\,0\,\% & \cumple{parcial}\,25\,\% & Más difícil de calcular, y algunos tramos pueden quedar tapado en la fotografía. \\
    \midrule
    \multicolumn{1}{l}{\sffamily\footnotesize\bfseries\color{uoc-azul}Anchura} & \textbf{45\,\%} & \textbf{95\,\%} & \\
    \glslink{ancho-libre-paso}{Anchura libre de paso} ($\geq$~1,80\,m) & \cumple{si}\,90\,\% & \cumple{si}\,90\,\% & Es una medida con facilidad de obtención en ambos tipos de imagen, aunque árboles o toldos pueden taparla. \\
    \arrayrulecolor{black!25}\cmidrule(lr){1-4}\arrayrulecolor{black}
    \Gls{altura-libre-paso} ($\geq$~2,20\,m) & \cumple{no}\,0\,\% & \cumple{si}\,100\,\% & Solo se aprecia a nivel de calle. \\
    \midrule
    \multicolumn{1}{l}{\sffamily\footnotesize\bfseries\color{uoc-azul}Pavimento} & \textbf{40\,\%} & \textbf{30\,\%} & \\
    Resbaladicidad (clase 3 del CTE) & \cumple{no}\,0\,\% & \cumple{no}\,0\,\% & No es posible extraerla. \\
    \arrayrulecolor{black!25}\cmidrule(lr){1-4}\arrayrulecolor{black}
    Continuidad y resaltes (hasta 4\,mm) & \cumple{no}\,0\,\% & \cumple{parcial}\,25\,\% & Solo se aprecia si la fotografía tiene una resolución extremadamente alta, porque 4\,mm es una medida muy precisa. \\
    \arrayrulecolor{black!25}\cmidrule(lr){1-4}\arrayrulecolor{black}
    \glslink{pavimento-tactil}{Pavimento táctil} de encaminamiento (4\,mm de altura, relieve lineal) & \cumple{parcial}\,75\,\% & \cumple{parcial}\,50\,\% & Desde el aire se ve la forma, pero no los milímetros. A pie de calle, 4\,mm es una medida muy precisa y es difícil de distinguir todo si no hay varias imágenes. \\
    \arrayrulecolor{black!25}\cmidrule(lr){1-4}\arrayrulecolor{black}
    Pavimento táctil de advertencia (4\,mm de altura, botones o franjas) & \cumple{parcial}\,75\,\% & \cumple{parcial}\,50\,\% & Desde el aire se ven potencialmente los botones y franjas, no los milímetros. A pie de calle, 4\,mm es una medida muy precisa. \\
    \midrule
    \multicolumn{1}{l}{\sffamily\footnotesize\bfseries\color{uoc-azul}\glslink{cruce-peatonal}{Cruces}} & \textbf{70\,\%} & \textbf{50\,\%} & \\
    Dimensiones del vado (anchura mínima de 1,80\,m, resalte inferior a 4\,mm) & \cumple{parcial}\,75\,\% & \cumple{parcial}\,25\,\% & Desde el aire se ve la anchura, pero no el resalte; a pie de calle se ve el resalte, pero puede quedar tapado fácilmente. \\
    \arrayrulecolor{black!25}\cmidrule(lr){1-4}\arrayrulecolor{black}
    Dimensiones del cruce (anchura igual o mayor que los vados, mínimo 1,80\,m) & \cumple{si}\,100\,\% & \cumple{parcial}\,25\,\% & Desde el aire se ve la anchura completa; a pie de calle haría falta una imagen lateral o un cálculo más complicado. \\
    \arrayrulecolor{black!25}\cmidrule(lr){1-4}\arrayrulecolor{black}
    Semaforización accesible (pulsador entre 0,80 y 1,20\,m, avisador acústico) & \cumple{no}\,0\,\% & \cumple{parcial}\,50\,\% & A pie de calle solo se podría estimar la altura del pulsador, no el avisador acústico. \\
    \arrayrulecolor{black!25}\cmidrule(lr){1-4}\arrayrulecolor{black}
    \glslink{isleta-refugio}{Isletas de refugio} (obligatorias si el cruce supera 14\,m) & \cumple{si}\,100\,\% & \cumple{si}\,100\,\% & Se ven desde ambos tipos de imagen. \\
    \midrule
    \multicolumn{1}{l}{\sffamily\footnotesize\bfseries\color{uoc-azul}\glslink{mobiliario-urbano}{Mobiliario}} & \textbf{65\,\%} & \textbf{80\,\%} & \\
    Ubicación del mobiliario (alineado, fuera del itinerario peatonal accesible) & \cumple{si}\,90\,\% & \cumple{si}\,80\,\% & Se ve en planta; los árboles pueden taparlo. \\
    \arrayrulecolor{black!25}\cmidrule(lr){1-4}\arrayrulecolor{black}
    Dimensiones y alturas (elementos elevados con altura libre $\geq$~2,20\,m) & \cumple{no}\,0\,\% & \cumple{parcial}\,75\,\% & A pie de calle se podría calcular, pero con dificultades. \\
    \arrayrulecolor{black!25}\cmidrule(lr){1-4}\arrayrulecolor{black}
    Franja libre de paso (anchura mínima del itinerario, 1,80\,m) & \cumple{si}\,90\,\% & \cumple{si}\,80\,\% & Se ve en planta, aunque los árboles pueden taparlo. \\
    \arrayrulecolor{black!25}\cmidrule(lr){1-4}\arrayrulecolor{black}
    Ocupación por terrazas (sin invadir el itinerario ni ocultar la señalización) & \cumple{si}\,90\,\% & \cumple{si}\,80\,\% & Se ve en planta, aunque los árboles pueden taparlo. \\
    \arrayrulecolor{black!25}\cmidrule(lr){1-4}\arrayrulecolor{black}
    Bolardos (alineados, altura entre 0,75 y 1,00\,m, color contrastado) & \cumple{parcial}\,50\,\% & \cumple{parcial}\,75\,\% & Desde el aire se puede calcular su alineación. A pie de calle se aprecia la altura, pero es más difícil calcular la alineación. \\
    \midrule
    \multicolumn{1}{l}{\sffamily\footnotesize\bfseries\color{uoc-azul}\glslink{senalizacion}{Señalización}} & \textbf{10\,\%} & \textbf{50\,\%} & \\
    Legibilidad, contraste y tamaño según la distancia (0,7 a 7,0\,cm) & \cumple{no}\,0\,\% & \cumple{parcial}\,75\,\% & A pie de calle, se podría llegar a estimar si la fotografía tiene resolución suficiente. \\
    \arrayrulecolor{black!25}\cmidrule(lr){1-4}\arrayrulecolor{black}
    Altura de los paneles (entre 0,90 y 1,20\,m) & \cumple{no}\,0\,\% & \cumple{si}\,100\,\% & A pie de calle se puede si la fotografía tiene resolución suficiente. \\
    \arrayrulecolor{black!25}\cmidrule(lr){1-4}\arrayrulecolor{black}
    Información escrita (lectura fácil, pictogramas, nombres de las vías en los cruces) & \cumple{no}\,0\,\% & \cumple{parcial}\,75\,\% & Necesita fotografías de alta resolución. \\
    \arrayrulecolor{black!25}\cmidrule(lr){1-4}\arrayrulecolor{black}
    Señalización táctil (braille y alto relieve) & \cumple{no}\,0\,\% & \cumple{parcial}\,10\,\% & El relieve apenas es visible en una imagen. Se necesita inspección in situ. \\
    \arrayrulecolor{black!25}\cmidrule(lr){1-4}\arrayrulecolor{black}
    Información sonora (señalización acústica en cruces) & \cumple{no}\,0\,\% & \cumple{no}\,0\,\% & No es observable en una imagen. \\
    \arrayrulecolor{black!25}\cmidrule(lr){1-4}\arrayrulecolor{black}
    Elementos transparentes (bandas opacas de 5 a 10\,cm) & \cumple{parcial}\,50\,\% & \cumple{parcial}\,50\,\% & Potencialmente posible, pero difícil. \\
    \midrule
    \multicolumn{1}{l}{\sffamily\footnotesize\bfseries\color{uoc-azul}Iluminación} & \textbf{15\,\%} & \textbf{30\,\%} & \\
    Niveles de iluminación (RD 1890/2008) & \cumple{parcial}\,25\,\% & \cumple{parcial}\,50\,\% & Los fija el RD 1890/2008 \cite{rd_1890_2008,guia_mitma_2021}; solo serían estimables con imágenes nocturnas. \\
    \arrayrulecolor{black!25}\cmidrule(lr){1-4}\arrayrulecolor{black}
    Uniformidad luminosa & \cumple{parcial}\,25\,\% & \cumple{parcial}\,25\,\% & Solo estimable con imágenes nocturnas. \\
    \arrayrulecolor{black!25}\cmidrule(lr){1-4}\arrayrulecolor{black}
    Detección de obstáculos & \cumple{no}\,0\,\% & \cumple{parcial}\,10\,\% & Potencialmente, con imágenes de día y de noche. \\
\end{longtable}
\endgroup

Los artículos legales de cada requisito se encuentran en las tablas de la sección~\ref{sec:caracteristicas}, de la que solo se incluyen las características que se corresponden con las dimensiones puntuadas en los datos. Transporte, que solo existe en Albacete, no se incluye porque no se describe en esa sección. En la fila de cada dimensión se indica el porcentaje medio de sus requisitos, redondeado a múltiplos de 5. El símbolo resume el porcentaje: \cumple{no}~0\,\%, \cumple{parcial}~10--75\,\% y \cumple{si}~80--100\,\%.
\\

Según esta estimación, \textbf{ninguna dimensión se puede obtener por completo} a partir de imágenes, aunque algunas se quedan bastante cerca. Ningún requisito llega al 100\,\% desde ambos tipos de imagen, salvo las isletas de refugio, por los problemas que pueden presentar las imágenes, como la vegetación u otros obstáculos que oculten la escena (tabla~\ref{tab:viabilidad}).
\\

La dimensión más viable es la \textbf{anchura}, sobre todo a pie de calle (95\,\%; desde el aire, 45\,\%). El porcentaje aéreo es bajo porque incluye la altura libre de paso, que solo se aprecia a nivel de calle; la anchura libre de paso, por sí sola, alcanza el 90\,\% en ambos casos. Desde el aire, lo más alto son los cruces (70\,\%), el mobiliario (65\,\%) y la anchura (45\,\%). A pie de calle, tras la anchura, destacan el mobiliario (80\,\%), la señalización (50\,\%) y los cruces (50\,\%).
\\

Las dimensiones de menor viabilidad son la \textbf{iluminación} (15\,\% desde el aire y 30\,\% a pie de calle, y suponiendo que se dispone de imágenes nocturnas), la \textbf{pendiente} (0\,\% y 40\,\%) y el \textbf{pavimento} (40\,\% y 30\,\%). En el pavimento, la resbaladicidad no se puede extraer de ninguna manera, y los resaltes de 4\,mm solo en parte (25\,\% a pie de calle). La señalización se queda en el 10\,\% desde el aire y en el 50\,\% a pie de calle.
\\

Esta información servirá para determinar qué dimensiones se utilizan en las hipótesis de investigación (H1, H2, H3 y H5), comprobando la precisión del modelo usando dimensiones que se pueden capturar desde imágenes aéreas, a pie de calle o forzosamente in situ por el equipo técnico.
\\